% ============================================================
\section{Discussion}
\label{sec:discussion}
% ============================================================

\paragraph{What the results say.}
Quantum-inspired encodings, deployed as fixed feature maps with linear
classification, do not provide a statistically significant accuracy advantage
over classical feature engineering on any of ten diverse benchmark datasets.
The failure is not uniform: each encoding has a distinct, spectrally
identifiable cause.
The mechanistic attribution (rank collapse, geometric equivalence, Hamming
mismatch) is more informative than a simple accuracy ranking.

\paragraph{Qualified positive findings.}
Three findings deserve emphasis rather than dismissal.
First, encoding overhead is negligible and should not be cited as a reason to
avoid QIE in future work.
Second, amplitude encoding on high-rank noise is statistically
indistinguishable from the best classical method; this is a structurally
grounded result, not luck.
Third, basis encoding's distinctly low CKA with classical baselines suggests
it could serve as a complementary feature in ensemble settings rather than a
standalone replacement.

\paragraph{Implications for quantum ML.}
The results do not speak to quantum hardware advantage.
Actual quantum computers can exploit superposition and entanglement in ways
that are classically intractable.
What this work shows is that the \emph{geometric structure} induced by quantum
encoding circuits, in isolation, does not confer a representational advantage
on classical data.
If quantum advantage in ML materialises, it will not be because of the fixed
feature map alone.

\paragraph{Where quantum advantage may still appear.}
A natural reviewer objection is that the ten benchmark datasets used here
are not the tasks where quantum advantage would be expected to show up.
This objection is well-founded.
Tasks with genuinely quantum-structured inputs, such as molecular energy
surfaces, quantum simulation outputs, or problems with provably hard classical
representations, may exhibit a different picture.
The present findings are therefore best interpreted as ruling out fixed QIE
geometry as a source of advantage on \emph{classical} data distributions, not
as ruling out quantum advantage in general.
Whether quantum-structured data breaks the spectral attribution framework
developed here is an open and important question.

\paragraph{Limitations.}
This benchmark evaluates fixed, parameter-free feature maps with a linear
probe or shallow MLP.
Learned quantum-inspired encodings (parameterised quantum circuits or trainable
angle schedules) are out of scope and represent the most natural next step:
if the angle schedule is jointly optimised with the downstream loss, the
geometric-equivalence failure mode documented here may not apply.
All comparisons use a single downstream linear classifier or shallow MLP;
deeper architectures may interact differently with QIE representations.
Neural Tangent Kernel (NTK) analysis, which would characterise how each
encoding reshapes the optimisation landscape and whether observed conditioning
differences propagate to gradient-flow behaviour, was not conducted in this
study and remains an open direction for follow-on work.
